\chapter{Exceptional Yangians: quadratic RTT data and the duality problem}
\label{chap:exceptional-yangian-koszul-duality-platonic}

\providecommand{\ClaimStatusProvedHere}{\textsc{ProvedHere}}
\providecommand{\ClaimStatusProvedElsewhere}{\textsc{ProvedElsewhere}}
\providecommand{\ClaimStatusConditional}{\textsc{Conditional}}
\providecommand{\ClaimStatusComputed}{\textsc{Computed}}
\providecommand{\ClaimStatusDefinitional}{\textsc{Definitional}}
\providecommand{\PhiFA}{\Phi^{\mathrm{FA}}}
\providecommand{\SpCh}{\mathrm{Sp}^{\mathrm{ch}}}

Fix
\[
 \fg\in\{E_6,E_7,E_8,F_4,G_2\},
 \qquad Y_\hbar(\fg)\text{ over }\C[[\hbar]].
\]
The local calculations determine Cartan matrices, root counts, selected
Weyl dimensions, and four recorded tensor-square decompositions.  An
exceptional RTT theorem requires further structure: a type-specific
kernel, its transfer matrix, a presentation morphism to
$Y_\hbar(\fg)$, PBW and Koszul statements in that presentation, and a
calculation of the orthogonal relation space.  The central map is the
inverse-kernel comparison
$R_N(K_{\fg,\hbar})^\perp\to R_N(K_{\fg,\hbar}^{-1})$ and its passage
through the Yangian presentation map.

\begin{remark}[Affine and Hall companions; \ClaimStatusConditional]
\label{rem:exceptional-yangian-archetype-and-volIII-shadow}
\index{Yangian!affine companion}
\index{K3!Hall--Drinfeld companion}
The affine vertex algebra $V_k(\fg)$ belongs to the Open chart and has
the numerical row
\[
\bigl(\kappa_{\mathrm{cat}},\kappa^{\mathrm{Hodge}}_{\mathrm{ch}},
\kappa^{\mathrm{Heis}}_{\mathrm{ch}},\kappa_{\mathrm{BKM}},
\kappa_{\mathrm{fiber}}\bigr)
=\left(0,0,\frac{\dim\fg\,(k+h^\vee)}{2h^\vee},-,\operatorname{rk}\fg\right).
\]
The K3 Hall construction has the typed arrow
\[
\PhiFA_3\bigl(D^b\Coh(K3\times E)\bigr)
\xrightarrow{\ \SpCh_{\Sigma_2,C}\ }
\mathcal D_\hbar\!\left(
\mathcal Y^{\mathrm{Hall}}_\hbar(\mathrm{CoHA}_{K3\times E})
\right).
\]
The exceptional RTT problem is governed by a different arrow,
\[
\pi_{\fg}\colon
\widehat A^{\mathrm{RTT}}(K_{\fg,\hbar})/\mathcal N_{\fg}
\longrightarrow Y_\hbar(\fg).
\]
The RTT duality criterion uses $\pi_{\fg}$ and the quadratic relation
spaces of its source.  A comparison with the Hall construction consists
of a filtered bialgebra morphism into the displayed Drinfeld double.
\end{remark}

\section{Finite RTT presentations}
\label{sec:exceptional-yangian-finite-rtt-windows}

Let $V_0$ be a finite-dimensional vector space.  The homogeneous
extended RTT presentation uses
\[
W_N(V_0):=\End(V_0)\otimes
\operatorname{span}\{1,u^{-1},\ldots,u^{-N}\}.
\]
Writing
$T_N(u)=\sum_{r=0}^{N}T^{(r)}u^{-r}$, a kernel
\[
K(u)\in\End(V_0\otimes V_0)((u^{-1}))[[\hbar]],
\qquad K(u)=1+O(u^{-1}),
\]
defines a quadratic subspace $R_N(K)\subset W_N(V_0)^{\otimes2}$:
it is spanned by the coefficients, projected to modes at most $N$, of
\begin{equation}\label{eq:exceptional-rtt-relation-map}
K_{12}(u-v)T_{N,1}(u)T_{N,2}(v)
-T_{N,2}(v)T_{N,1}(u)K_{12}(u-v).
\end{equation}
Set
\[
A_N(K):=T(W_N(V_0))/\langle R_N(K)\rangle.
\]

\begin{definition}[Finite RTT presentation datum;
\ClaimStatusDefinitional]
\label{def:finite-yangian-module-packet}
\index{Yangian!finite RTT presentation datum}
A finite RTT presentation datum for $Y_\hbar(\fg)$ is a tuple
\[
\mathscr R_{\fg,N}
=(V_0,\rho_0,K_{\fg,\hbar}(u),T_{\fg,N}(u),
\mathcal N_{\fg,N},\pi_{\fg,N}).
\]
Here $\rho_0$ is a chosen finite-dimensional $\fg$-module,
$K_{\fg,\hbar}$ is a type-specific kernel on $V_0^{\otimes2}$,
$T_{\fg,N}$ is the universal transfer matrix in $A_N(K_{\fg,\hbar})$,
$\mathcal N_{\fg,N}$ is the stated normalization ideal, including the
specialization $T^{(0)}=I$ when this is part of the presentation, and
\[
\pi_{\fg,N}\colon
A_N(K_{\fg,\hbar})/\mathcal N_{\fg,N}
\xrightarrow{\ \sim\ }Y_{\hbar,\leq N}(\fg)
\]
is a filtered presentation isomorphism.  The transition maps
$p_{N+1,N}$ form part of the datum.  If
$r_{N+1,N}$ denotes the transition on the specified Yangian
truncations, they satisfy
$\pi_{\fg,N}p_{N+1,N}=r_{N+1,N}\pi_{\fg,N+1}$.
When the normalization introduces lower filtration terms, the datum
also contains the Rees algebra
\[
\mathcal R_{\fg,N}
:=\operatorname{Rees}_F\!\left(
A_N(K_{\fg,\hbar})/\mathcal N_{\fg,N}\right)
\]
and the specializations
\[
\mathcal R_{\fg,N}/(t-1)\xrightarrow{\sim}
Y_{\hbar,\le N}(\fg),
\qquad
\mathcal R_{\fg,N}/(t)\xrightarrow{\sim}
\operatorname{gr}_F Y_{\hbar,\le N}(\fg).
\]
\end{definition}

\begin{proposition}[Trace duality for the coefficient space;
\ClaimStatusProvedHere]
\label{prop:finite-rtt-trace-pairing-nondegenerate}
The formula
\[
\bigl\langle Xu^{-r},Yu^{-s}\bigr\rangle_N
=\operatorname{Tr}_{V_0}(XY)\,\delta_{rs}
\]
gives a perfect pairing on $W_N(V_0)$.  It therefore defines
\[
R_N(K)^\perp\subset W_N(V_0)^\vee\otimes W_N(V_0)^\vee.
\]
\end{proposition}

\begin{proof}
For matrix units $E_{ij}$ one has
$\operatorname{Tr}(E_{ij}E_{kl})=\delta_{jk}\delta_{il}$.
Thus $E_{ij}u^{-r}$ pairs with $E_{ji}u^{-r}$, and the Gram
matrix is a permutation matrix in every spectral mode.
\end{proof}

\begin{definition}[Spectral completion and inverse-kernel comparison;
\ClaimStatusDefinitional]
\label{def:finite-rtt-window-spectral-completion}
\index{RTT algebra!spectral completion}
The coefficientwise $(u^{-1},\hbar)$-adic completion is denoted
$\widehat A_N(K)$.  The inverse-kernel comparison at level $N$ is the
map of relation spaces
\begin{equation}\label{eq:exceptional-inverse-kernel-comparison}
\vartheta_{\fg,N}\colon
R_N(K_{\fg,\hbar})^\perp
\longrightarrow R_N(K_{\fg,\hbar}^{-1}).
\end{equation}
It is an isomorphism precisely when the two subspaces coincide under
the trace identification $W_N(V_0)^\vee\simeq W_N(V_0)$.
The scalar-gauge comparison is an identity
\begin{equation}\label{eq:exceptional-scalar-gauge}
K_{\fg,\hbar}(u)^{-1}
=f_{\fg}(u,\hbar)K_{\fg,-\hbar}(u),
\qquad
f_{\fg}\in\C((u^{-1}))[[\hbar]]^\times.
\end{equation}
\end{definition}

\begin{lemma}[Permutation-kernel benchmark; \ClaimStatusProvedHere]
Let $P(x\otimes y)=y\otimes x$ and
$K_\hbar(u)=1-\hbar P/u$.  Then
\[
K_\hbar(u)^{-1}
=\left(1-\frac{\hbar^2}{u^2}\right)^{-1}
\left(1+\frac{\hbar P}{u}\right)
=\left(1-\frac{\hbar^2}{u^2}\right)^{-1}K_{-\hbar}(u).
\]
\end{lemma}

\begin{proof}
The equality follows from $P^2=1$ by multiplying the two factors
$1-\hbar P/u$ and $1+\hbar P/u$.
\end{proof}

\section{Typed comparison maps}
\label{sec:exceptional-yangian-taxa}

\begin{proposition}[Objects entering the exceptional RTT problem;
\ClaimStatusProvedHere]
\label{prop:exceptional-yangian-type-separation}
The constructions used in this chapter have the following types:
\[
\begin{array}{c|c}
\text{object}&\text{ambient category}\\ \hline
(A_{\fg},\Phi_{\fg})&\text{based root data}\\
(V_0,\rho_0)&\operatorname{Rep}(\fg)\\
A_N(K_{\fg,\hbar})&\operatorname{Alg}^{\mathrm{quad}}_{\C[[\hbar]]}\\
Y_{\hbar,\leq N}(\fg)&\operatorname{HopfAlg}^{\mathrm{filt}}_{\C[[\hbar]]}\\
B(A_N)&\operatorname{Coalg}^{\mathrm{conil}}\\
V_k(\fg)&\operatorname{ChAlg}(X)\\
\mathcal D_\hbar(\mathcal Y^{\mathrm{Hall}}_\hbar(\mathrm{CoHA}_{K3\times E}))
&\operatorname{HopfAlg}^{\mathrm{comp}}.
\end{array}
\]
The arrows relating the first five rows are
\[
\rho_0,\qquad
\pi_{\fg,N}\colon A_N(K_{\fg,\hbar})/\mathcal N_{\fg,N}
\xrightarrow{\sim}Y_{\hbar,\leq N}(\fg),
\qquad
q_N\colon A_N^{\mathrm i}\longrightarrow B(A_N).
\]
A chiral conclusion additionally uses a strong monoidal
factorization functor $\mathfrak C_X$ and a comparison
\[
\gamma_N\colon
\mathfrak C_X\bigl(A_N^{!,\mathrm{quad}}\bigr)
\xrightarrow{\sim}
\mathfrak C_X(A_N)^{!,\mathrm{ch}}.
\]
\end{proposition}

\begin{proof}
Each row follows from its defining generators and structure maps.
The displayed arrows state the additional data required to move from
one row to the next.
\end{proof}

\section{The inverse-kernel criterion}
\label{sec:exceptional-yangian-template}

\begin{proposition}[Quadratic RTT duality criterion;
\ClaimStatusProvedHere]
\label{prop:exceptional-yangian-template}
This is the algebraic RTT presentation of the Open chart at levels
$1\leftrightarrow2$.  The finite hypothesis package is
$H_{\mathrm{RTT},N}=\textup{(a)}$--$\textup{(d)}$; the completed
package adds transition compatibility and Mittag--Leffler, and the
chiral package adds $(\mathfrak C_X,\gamma_N)$.
Let $\mathscr R_{\fg,N}$ be a finite RTT presentation datum.  Assume:
\begin{enumerate}[label=\textup{(\alph*)}]
\item $\pi_{\fg,N}$ is a filtered algebra isomorphism, and the
normalization ideal is compatible with quadratic duality and with
$\hbar\mapsto-\hbar$;
\item the same presentation has a PBW isomorphism
\[
\operatorname{gr}_F Y_{\hbar,\le N}(\fg)
\xrightarrow{\sim}U(\fg[t])_{\le N}
\]
and diagonal Koszul cohomology
\[
H^{p,q}B(A_N)=0\qquad(p\ne q),
\]
so that $q_N\colon A_N^{\mathrm i}\to B(A_N)$ is a quasi-isomorphism;
\item $\vartheta_{\fg,N}$ in
\eqref{eq:exceptional-inverse-kernel-comparison} is an isomorphism;
\item the scalar-gauge identity
\eqref{eq:exceptional-scalar-gauge} holds.
\end{enumerate}
Then the quadratic dual algebra satisfies
\begin{equation}\label{eq:exceptional-finite-quadratic-dual}
A_N(K_{\fg,\hbar})^{!,\mathrm{quad}}
\cong A_N(K_{\fg,-\hbar}).
\end{equation}
The duality-compatible normalization and $\pi_{\fg,N}$ give
\[
Y_{\hbar,\le N}(\fg)^{!,\mathrm{quad}}
\cong Y_{-\hbar,\le N}(\fg).
\]
If the isomorphisms commute with $p_{N+1,N}$ and every finite total
weight window is Mittag--Leffler, they assemble to
\begin{equation}\label{eq:exceptional-completed-quadratic-dual}
\widehat A(K_{\fg,\hbar})^{!,\mathrm{cont}}
\cong\widehat A(K_{\fg,-\hbar}).
\end{equation}
If $\mathfrak C_X$ and $\gamma_N$ of
Proposition~\ref{prop:exceptional-yangian-type-separation} are supplied,
the same isomorphism holds in the stated $E_1$-chiral completion.
\end{proposition}

\begin{proof}
For $A_N=T(W_N)/\langle R_N\rangle$, quadratic duality gives
\[
A_N^{!,\mathrm{quad}}
=T(W_N^\vee)/\langle R_N^\perp\rangle.
\]
Hypothesis~\textup{(c)} replaces $R_N^\perp$ by the relation space of
$K_{\fg,\hbar}^{-1}$; hypothesis~\textup{(d)} replaces this kernel by
$K_{\fg,-\hbar}$ because a central scalar factor cancels from
\eqref{eq:exceptional-rtt-relation-map}.  Hypothesis~\textup{(b)}
identifies the quadratic coalgebra with bar cohomology.
Hypothesis~\textup{(a)} carries the result through the normalization
and the presentation isomorphism.  Compatibility
with $p_{N+1,N}$ and the Mittag--Leffler condition gives the completed
isomorphism weight by weight.  The comparison $\gamma_N$ transports it
to the chiral category.
\end{proof}

\begin{remark}[Chevalley involution]
\label{rem:exceptional-chevalley-intrinsic}
The Lie-algebra involution
\[
e_i\longmapsto-f_i,\qquad f_i\longmapsto-e_i,
\qquad h_i\longmapsto-h_i
\]
acts by $-I$ on the simple-root lattice.  Its use in
Proposition~\ref{prop:exceptional-yangian-template} requires a compatible
Yangian morphism
\[
\omega_{\fg,N}\colon
A_N(K_{\fg,\hbar})\longrightarrow
A_N(K_{\fg,-\hbar})^{\mathrm{op}}
\]
intertwining $T_{\fg,N}$, $\mathcal N_{\fg,N}$, and the transition maps.
This is a presentation-level construction.
\end{remark}

\begin{lemma}[Passage from finite windows to the completed algebra;
\ClaimStatusProvedHere]
\label{lem:exceptional-finite-window-promotion}
Let $\theta_N\colon A_N^{!,\mathrm{quad}}\to A_N'$ be finite-window
isomorphisms commuting with the transition maps.  Suppose each finite
total weight $w$ is eventually constant in $N$ on both sides.  Then
\[
\widehat\theta\colon
\varprojlim_N A_N^{!,\mathrm{quad}}
\xrightarrow{\sim}\varprojlim_N A_N'
\]
is an isomorphism in the completed graded category.
\end{lemma}

\begin{proof}
For fixed $w$, choose $N\ge w$.  Both inverse systems are constant in
that weight, and $\theta_N$ identifies their stable values.  Taking the
product over $w$ gives $\widehat\theta$.  Equivalently, the Milnor
sequence has $\varprojlim^{1}=0$ in every weight.
\end{proof}

\begin{proposition}[Ordered bar information under averaging;
\ClaimStatusProvedHere]
\label{prop:exceptional-yangian-ordered-bar-averaging}
For $A_N=T(W_N)/\langle R_N\rangle$, the ordered reduced bar coalgebra is
\[
B^{\mathrm{ord}}(A_N)=
\bigl(T^c(s^{-1}\bar A_N),d_\mu,\Delta_{\mathrm{dec}}\bigr).
\]
The quotient
\[
\operatorname{Av}_N\colon
T^c(s^{-1}\bar A_N)\longrightarrow
\bigoplus_{m\ge0}(s^{-1}\bar A_N)^{\otimes m}_{\Sigma_m}
\]
has kernel containing
$x\otimes y-(-1)^{|x||y|}y\otimes x$.  Hence the equality
$R_N(K)^\perp=R_N(K^{-1})$ is an equality in the ordered tensor square.
An averaged formulation requires a section of $\operatorname{Av}_N$ or
an equivariant descent theorem carrying both relation spaces.
\end{proposition}

\begin{proof}
Deconcatenation records the order of the two tensor factors, whereas
coinvariants identify their graded permutations.  The orthogonal
relation space belongs to $W_N^\vee\otimes W_N^\vee$, so its recovery
from coinvariants is exactly the stated section or descent datum.
\end{proof}

\section{Root data and the scope of the cited PBW source}
\label{sec:exceptional-yangian-pbw}

\begin{proposition}[Guay--Regelskis--Wendlandt;
\ClaimStatusProvedElsewhere]
\label{thm:exceptional-yangian-pbw-grw18}
Guay--Regelskis--Wendlandt~\cite{GRW18} compare the RTT, extended
Drinfeld, and Drinfeld presentations for orthogonal and symplectic
Yangians.  Their Lie types are $B$, $C$, and $D$.  For an exceptional
$\fg$, the corresponding input to
Proposition~\ref{prop:exceptional-yangian-template} is an isomorphism
\[
\pi_{\fg,N}\colon
A_N(K_{\fg,\hbar})/\mathcal N_{\fg,N}
\xrightarrow{\sim}Y_{\hbar,\le N}(\fg)
\]
together with PBW in this same presentation.
\end{proposition}

\begin{corollary}[Exceptional Cartan calculations;
\ClaimStatusComputed]
\label{cor:exceptional-yangian-graded}
\index{exceptional root systems!Cartan matrices}
In the numbering used by
\texttt{compute/lib/yangian\_rtt\_exceptional.py}, the stored Cartan
matrices are
\[
A_{G_2}=\begin{pmatrix}2&-1\\-3&2\end{pmatrix},\qquad
A_{F_4}=\begin{pmatrix}
2&-1&0&0\\-1&2&-2&0\\0&-1&2&-1\\0&0&-1&2
\end{pmatrix},
\]
\[
A_{E_6}=\begin{pmatrix}
2&0&-1&0&0&0\\0&2&0&-1&0&0\\-1&0&2&-1&0&0\\
0&-1&-1&2&-1&0\\0&0&0&-1&2&-1\\0&0&0&0&-1&2
\end{pmatrix},
\]
\[
A_{E_7}=\begin{pmatrix}
2&0&-1&0&0&0&0\\0&2&0&-1&0&0&0\\-1&0&2&-1&0&0&0\\
0&-1&-1&2&-1&0&0\\0&0&0&-1&2&-1&0\\
0&0&0&0&-1&2&-1\\0&0&0&0&0&-1&2
\end{pmatrix},
\]
\[
A_{E_8}=\begin{pmatrix}
2&0&-1&0&0&0&0&0\\0&2&0&-1&0&0&0&0\\
-1&0&2&-1&0&0&0&0\\0&-1&-1&2&-1&0&0&0\\
0&0&0&-1&2&-1&0&0\\0&0&0&0&-1&2&-1&0\\
0&0&0&0&0&-1&2&-1\\0&0&0&0&0&0&-1&2
\end{pmatrix}.
\]
Iterated simple reflections in the program give
\[
\begin{array}{c|ccccc}
\fg&E_6&E_7&E_8&F_4&G_2\\ \hline
\operatorname{rk}\fg&6&7&8&4&2\\
|\Phi_+|&36&63&120&24&6\\
\dim\fg=\operatorname{rk}\fg+2|\Phi_+|&78&133&248&52&14\\
d& (1,1,1,1,1,1)&(1,1,1,1,1,1,1)&(1,1,1,1,1,1,1,1)&(1,1,2,2)&(3,1)\\
|\operatorname{Out}(\fg)|&2&1&1&1&1
\end{array}
\]
The reference table in the same module stores
\[
(h^\vee_{E_6},h^\vee_{E_7},h^\vee_{E_8},h^\vee_{F_4},h^\vee_{G_2})
=(12,18,30,9,4).
\]
Consequently
\[
\kappa(V_k(\fg))=\frac{\dim\fg\,(k+h^\vee)}{2h^\vee}
\]
gives
\[
\frac{13(k+12)}4,\quad
\frac{133(k+18)}{36},\quad
\frac{62(k+30)}{15},\quad
\frac{26(k+9)}9,\quad
\frac{7(k+4)}4
\]
for $E_6,E_7,E_8,F_4,G_2$, respectively.
\end{corollary}

\begin{proof}
The program enumerates $\Phi_+$ from the displayed matrices and checks
$\dim\fg=\operatorname{rk}\fg+2|\Phi_+|$.  It finds the diagonal
symmetrizers $d$ from $d_iA_{ij}=d_jA_{ji}$ and counts permutations
preserving each Cartan matrix.  The values of $h^\vee$ are reference
inputs in \texttt{EXCEPTIONAL\_DATA}.  Substitution in the formula for
$\kappa$ gives the final row.
\end{proof}

\begin{proposition}[Representation calculations in the local programs;
\ClaimStatusComputed]
\label{prop:exceptional-fundamental-packet-scope}
\index{exceptional Yangian!representation calculations}
The Weyl-dimension routines evaluate
\[
\dim V_{E_6}(\omega_1)=\dim V_{E_6}(\omega_6)=27,
\quad
\dim V_{E_7}(\omega_7)=56,
\quad
\dim V_{E_8}(\omega_8)=248.
\]
The rank-two routine gives $\dim V_{G_2}(\omega_1)=7$.
The program then instantiates the following recorded decompositions and
checks the dimension sums and the displayed Casimir eigenvalues:
\[
\begin{array}{c|c|c}
\fg&V_0^{\otimes2}&C_2\text{ on the listed summands}\\ \hline
E_6&351_A\oplus351_S\oplus27^*_S&100/3,\ 112/3,\ 52/3\\
E_7&1_A\oplus1539_A\oplus133_S\oplus1463_S&0,\ 56,\ 36,\ 60\\
E_8&1_S\oplus3875_S\oplus27000_S\oplus248_A\oplus30380_A
&0,\ 96,\ 124,\ 60,\ 120\\
G_2&1_S\oplus27_S\oplus7_A\oplus14_A&0,\ 28/3,\ 4,\ 8.
\end{array}
\]
For $F_4$, the function
\texttt{tensor\_product\_decomposition("F4")} returns the empty list.
Thus the first $F_4$ construction is a decomposition of
$V_{F_4}(\omega_4)^{\otimes2}$ together with explicit primitive
idempotents in $\End(V_{F_4}(\omega_4)^{\otimes2})$.
\end{proposition}

\begin{proof}
The highest weights, dimensions, symmetry tags, and Casimir values are
read directly from
\texttt{compute/lib/exceptional\_yangian\_engine.py}.  For the four
displayed types the code evaluates Weyl dimensions and
$C_2(\lambda)=(\lambda,\lambda+2\rho)$ on a decomposition supplied as
input.  It constructs dimensions and scalar eigenvalues; matrix
projectors $P_\lambda$, transfer matrices $T(u)$, and exceptional
kernels $K_{\fg,\hbar}(u)$ form the next constructions.
\end{proof}

\section{The five exceptional types}
\label{sec:exceptional-yangian-per-type}

\begin{proposition}[$E_6$ reduction; \ClaimStatusConditional]
\label{prop:exceptional-yangian-koszul-E6}
For $V_0=V(\omega_1)=27$ and $V_0^\vee=V(\omega_6)=27^*$, construct
\[
K_{E_6,\hbar}(u),\quad T_{E_6,N}(u),\quad
\pi_{E_6,N},\quad\vartheta_{E_6,N},\quad f_{E_6}(u,\hbar)
\]
on the three channels
$351_A\oplus351_S\oplus27^*_S$ and their contragredient channels.
Together with PBW, diagonal Koszulness, and compatible transition maps,
these data give
\[
\widehat Y_\hbar(E_6)^{!,\mathrm{cont}}
\cong\widehat Y_{-\hbar}(E_6).
\]
\end{proposition}

\begin{proof}
Apply Proposition~\ref{prop:exceptional-yangian-template} to the stated
three-channel presentation.
\end{proof}

\begin{remark}[$E_6$ contragredient representation]
\label{rem:exceptional-yangian-E6-outer}
The diagram involution exchanges $V(\omega_1)$ and $V(\omega_6)$.
The inverse-kernel map $\vartheta_{E_6,N}$ therefore carries an explicit
identification between the $27$ presentation and its contragredient
$27^*$ presentation, in addition to the diagram involution.
\end{remark}

\begin{proposition}[$E_7$ reduction; \ClaimStatusConditional]
\label{prop:exceptional-yangian-koszul-E7}
For the symplectic representation $V_0=V(\omega_7)=56$, construct the
four primitive projectors
\[
P_1,\ P_{1539},\ P_{133},\ P_{1463},
\qquad
\sum_\lambda P_\lambda=1,
\quad P_\lambda P_\mu=\delta_{\lambda\mu}P_\lambda,
\]
and a kernel
$K_{E_7,\hbar}(u)=\sum_\lambda k_\lambda(u,\hbar)P_\lambda$.
The maps $T_{E_7,N}$, $\pi_{E_7,N}$, $\vartheta_{E_7,N}$, together
with PBW, diagonal Koszulness, scalar gauge, and tower compatibility,
give
\[
\widehat Y_\hbar(E_7)^{!,\mathrm{cont}}
\cong\widehat Y_{-\hbar}(E_7).
\]
\end{proposition}

\begin{proof}
The four idempotents define $R_N(K)$ channel by channel; the conclusion
is Proposition~\ref{prop:exceptional-yangian-template}.
\end{proof}

\begin{remark}[$E_7$ symplectic form]
\label{rem:exceptional-yangian-E7-self-dual}
The invariant alternating form identifies $V_0$ with $V_0^\vee$.
A spectral kernel on this decomposition has four coefficient functions
$k_\lambda(u,\hbar)$; its inverse is governed by their reciprocals and
the four projector matrices.
\end{remark}

\begin{proposition}[$E_8$ reduction; \ClaimStatusConditional]
\label{prop:exceptional-yangian-koszul-E8}
For $V_0=\fg=248$, construct the five projectors on
\[
248\otimes248
=1_S\oplus3875_S\oplus27000_S\oplus248_A\oplus30380_A
\]
and a five-channel kernel $K_{E_8,\hbar}$.  A presentation map
$\pi_{E_8,N}$, PBW, diagonal Koszulness, the relation-space map
$\vartheta_{E_8,N}$, scalar gauge, and compatible finite windows give
\[
\widehat Y_\hbar(E_8)^{!,\mathrm{cont}}
\cong\widehat Y_{-\hbar}(E_8).
\]
\end{proposition}

\begin{proof}
This is the five-channel instance of
Proposition~\ref{prop:exceptional-yangian-template}.
\end{proof}

\begin{remark}[$E_8$ adjoint tensor square]
\label{rem:exceptional-yangian-E8-adjoint}
The equality
$1+3875+27000+248+30380=248^2=61504$
checks the dimensions of the five isotypic summands.  The RTT relation
space is obtained from their actual projector matrices and the five
spectral functions $k_\lambda(u,\hbar)$.
\end{remark}

\begin{proposition}[$F_4$ reduction; \ClaimStatusConditional]
\label{prop:exceptional-yangian-koszul-F4}
Let $V_0=V(\omega_4)$ have dimension $26$.  Construct its tensor-square
decomposition, the corresponding primitive projectors, and a spectral
kernel $K_{F_4,\hbar}$.  A presentation map $\pi_{F_4,N}$, PBW,
diagonal Koszulness, $\vartheta_{F_4,N}$, scalar gauge, and compatible
finite windows then give
\[
\widehat Y_\hbar(F_4)^{!,\mathrm{cont}}
\cong\widehat Y_{-\hbar}(F_4).
\]
\end{proposition}

\begin{proof}
The root datum supplies $A_{F_4}$, $d=(1,1,2,2)$,
$\dim F_4=52$, and $h^\vee=9$.  The listed representation and RTT
constructions supply hypotheses~\textup{(a)}--\textup{(d)} of
Proposition~\ref{prop:exceptional-yangian-template}.
\end{proof}

\begin{remark}[$F_4$ folding]
\label{rem:exceptional-yangian-F4-folding}
The fixed Lie algebra $E_6^{\mathbb Z/2}\cong F_4$ suggests a second
route.  That route consists of transition-compatible isomorphisms
\[
A_N(K_{E_6,\hbar})^{\mathbb Z/2}
\xrightarrow{\sim}A_N(K_{F_4,\hbar}),
\qquad
R_N(K_{E_6,\hbar})^{\perp,\mathbb Z/2}
\xrightarrow{\sim}R_N(K_{F_4,\hbar})^\perp.
\]
Exactness of invariants and compatibility with $\vartheta_N$ complete
the folding construction.
\end{remark}

\begin{proposition}[$G_2$ reduction; \ClaimStatusConditional]
\label{prop:exceptional-yangian-koszul-G2}
For $V_0=V(\omega_1)=7$, construct the four projectors on
\[
7\otimes7=1_S\oplus27_S\oplus7_A\oplus14_A
\]
and a four-channel kernel $K_{G_2,\hbar}$.  A presentation map
$\pi_{G_2,N}$, PBW, diagonal Koszulness, $\vartheta_{G_2,N}$,
scalar gauge, and compatible finite windows give
\[
\widehat Y_\hbar(G_2)^{!,\mathrm{cont}}
\cong\widehat Y_{-\hbar}(G_2).
\]
\end{proposition}

\begin{proof}
Apply Proposition~\ref{prop:exceptional-yangian-template} to the four
isotypic channels.
\end{proof}

\begin{remark}[$G_2$ triality]
\label{rem:exceptional-yangian-G2-folding}
The fixed Lie algebra $D_4^{\mathbb Z/3}\cong G_2$ gives a folding route
through maps
\[
A_N(K_{D_4,\hbar})^{\mathbb Z/3}
\longrightarrow A_N(K_{G_2,\hbar}).
\]
An isomorphism here, exactness of invariants, and compatibility with
$R_N(-)^\perp$ and $p_{N+1,N}$ provide the required descent theorem.
\end{remark}

\section{Exceptional-family reduction theorem}
\label{sec:exceptional-yangian-five-family}

\begin{theorem}[Exceptional Yangian duality from RTT data;
\ClaimStatusConditional]
\label{thm:exceptional-yangian-koszul-duality-all-five-types}
\index{Yangian!exceptional quadratic duality}
The source is the Open chart in algebraic RTT presentation, the target
is its continuous quadratic dual at levels $1\leftrightarrow2$, and the
hypothesis package is
$H_{\mathrm{excRTT}}=\textup{(1)}$--$\textup{(7)}$.
For each
$\fg\in\{E_6,E_7,E_8,F_4,G_2\}$, suppose the following constructions
are given in one presentation:
\begin{enumerate}[label=\textup{(\arabic*)}]
\item the finite RTT datum $\mathscr R_{\fg,N}$, the presentation
isomorphism $\pi_{\fg,N}$, and compatibility of
$\mathcal N_{\fg,N}$ with quadratic duality and
$\hbar\mapsto-\hbar$;
\item primitive projector matrices for every irreducible channel of
$V_0^{\otimes2}$ and a type-specific kernel $K_{\fg,\hbar}$ satisfying
the Yang--Baxter equation;
\item PBW and diagonal Koszulness for $A_N(K_{\fg,\hbar})$;
\item the relation-space isomorphism
$\vartheta_{\fg,N}\colon R_N(K_{\fg,\hbar})^\perp
\xrightarrow{\sim}R_N(K_{\fg,\hbar}^{-1})$;
\item the scalar-gauge identity
$K_{\fg,\hbar}^{-1}=f_{\fg}K_{\fg,-\hbar}$;
\item compatibility with $p_{N+1,N}$ and the Mittag--Leffler property in
each finite total weight;
\item for the chiral formulation, the factorization functor
$\mathfrak C_X$ and the comparison isomorphisms $\gamma_N$.
\end{enumerate}
Then
\[
\widehat Y_\hbar(\fg)^{!,\mathrm{cont}}
\cong\widehat Y_{-\hbar}(\fg),
\]
and the same formula holds in the specified $E_1$-chiral completion
under construction~\textup{(7)}.
\end{theorem}

\begin{proof}
The five propositions of
\S\ref{sec:exceptional-yangian-per-type} identify $V_0$ and the
isotypic channels.  Constructions~\textup{(1)}--\textup{(5)} give
Proposition~\ref{prop:exceptional-yangian-template} at every $N$;
construction~\textup{(6)} gives
Lemma~\ref{lem:exceptional-finite-window-promotion};
construction~\textup{(7)} transports the completed algebraic result to
the chiral category.
\end{proof}

\begin{corollary}[Passage from the classical families to all simple
types; \ClaimStatusConditional]
\label{cor:exceptional-yangian-all-simple}
\index{Yangian!all simple Lie types}
Molev~\cite{Molev07} gives classical RTT structures, and
Guay--Regelskis--Wendlandt~\cite{GRW18} compare three presentations
in types $B$, $C$, and $D$.  Assume the quadratic relation-space,
Koszul, and completion hypotheses of
Proposition~\ref{prop:exceptional-yangian-template} throughout the
classical families.  Adding the seven constructions of
Theorem~\ref{thm:exceptional-yangian-koszul-duality-all-five-types} for
$E_6,E_7,E_8,F_4,G_2$ yields the corresponding statement for every
simple complex Lie algebra.
\end{corollary}

\begin{remark}[Adjacent representation-theoretic constructions]
\label{rem:exceptional-yangian-does-not-close}
Shifted Yangians, BFN Coulomb branches, quantum affine algebras,
toroidal algebras, and elliptic Yangians carry their own presentations
and completions.  Comparisons with the theorem above are functors from
those presentations to the finite RTT systems $\mathscr R_{\fg,N}$.
\end{remark}

\begin{proposition}[Comparison data supplied by neighboring theories;
\ClaimStatusConditional]
\label{prop:exceptional-yangian-no-dk-root-k3-promotion}
A transfer from a neighboring construction to the exceptional RTT
theorem is a functor producing the tuple
\[
\mathscr R_{\fg,N}
=(V_0,\rho_0,K_{\fg,\hbar},T_{\fg,N},
\mathcal N_{\fg,N},\pi_{\fg,N})
\]
and the relation-space map
$\vartheta_{\fg,N}$.  The available source objects are:
\begin{enumerate}[label=\textup{(\alph*)}]
\item Drinfeld--Kohno supplies monodromy and braiding on tensor powers
of evaluation modules.
\item A root-of-unity specialization supplies an integral form and its
specialized braided category.
\item The K3 construction supplies the completed Hall--Drinfeld
bialgebra
\[
\mathcal D_\hbar\bigl(
\mathcal Y^{\mathrm{Hall}}_\hbar(\mathrm{CoHA}_{K3\times E})
\bigr).
\]
\item Affine, toroidal, elliptic, and Borcherds constructions supply
their respective filtered algebras, shuffle products, and denominator
identities.
\end{enumerate}
For each source, the comparison theorem is the construction of
$K_{\fg,\hbar}$, its projector matrices, $\pi_{\fg,N}$, and
$\vartheta_{\fg,N}$ with the transition identities of
Definition~\ref{def:finite-yangian-module-packet}.
\end{proposition}

\begin{proof}
The target tuple is exactly the input of
Proposition~\ref{prop:exceptional-yangian-template}; preservation of its
filtration and relation spaces gives the required hypothesis package.
\end{proof}

\section{Local calculations}
\label{sec:exceptional-yangian-verification}

The two local modules perform the following computations.
\begin{enumerate}[label=\textup{(\alph*)}]
\item \texttt{ExceptionalRootSystem} enumerates positive roots from the
five stored Cartan matrices and computes their symmetrizers.
\item \texttt{weyl\_dim\_explicit} evaluates the Weyl product for the
displayed $E_6$, $E_7$, and $E_8$ highest weights; the rank-two code
evaluates the $G_2$ dimensions.
\item \texttt{tensor\_product\_decomposition} stores the $E_6$, $E_7$,
$E_8$, and $G_2$ decompositions, evaluates component dimensions and
Casimir scalars, and returns an empty list for $F_4$.
\item The permutation-kernel test evaluates
$(1-\hbar P/u)(1+\hbar P/u)=(1-\hbar^2/u^2)I$.
\item The test file compares the root counts, dimensions, and diagram
automorphism orders with a stored reference dictionary and checks the
involution $(-I)^2=I$.
\item The routine named PBW consistency evaluates
\texttt{rs.dim\_algebra == GROUND\_TRUTH[name][dim]}; its output is the
classical dimension comparison in
Corollary~\ref{cor:exceptional-yangian-graded}.
\item
\texttt{SpectralDecomposition.yang\_baxter\_spectral\_check} assigns
\texttt{passes} from
$\sum_\lambda\dim V_\lambda=(\dim V_0)^2$.  The Yang--Baxter equation
appears separately as construction~\textup{(2)} in
Theorem~\ref{thm:exceptional-yangian-koszul-duality-all-five-types}.
\end{enumerate}

\begin{remark}[What the local tests compute]
\label{rem:exceptional-yangian-hziv}
The assertions in
\texttt{compute/tests/test\_exceptional\_yangian\_koszul\_duality.py}
establish the arithmetic listed above.  The exceptional theorem has
status \ClaimStatusConditional.  Its additional inputs are the matrices
$P_\lambda$, kernels $K_{\fg,\hbar}$, transfer matrices $T_{\fg,N}$,
presentation maps $\pi_{\fg,N}$, relation maps $\vartheta_{\fg,N}$,
bar differentials, and PBW filtrations.
\end{remark}

\section{Exceptional RTT data}
\label{sec:exceptional-yangian-summary}

\begin{table}[ht]
\centering
\small
\caption{Computed representation data and remaining RTT constructions.}
\label{tab:exceptional-yangian-theorem-status}
\begin{tabular}{c|c|p{4.1cm}|p{5.2cm}}
$\fg$&$\dim V_0$&\textsc{local calculation}&\textsc{remaining construction}\\ \hline
$E_6$&$27,27^*$&three recorded channels and their Casimir scalars&
projector matrices, $K,T,\pi$, PBW/Koszul, $\vartheta$, tower maps\\
$E_7$&$56$&four recorded channels and their Casimir scalars&
four projectors, $K,T,\pi$, PBW/Koszul, $\vartheta$, tower maps\\
$E_8$&$248$&five recorded channels and their Casimir scalars&
five projectors, $K,T,\pi$, PBW/Koszul, $\vartheta$, tower maps\\
$F_4$&$26$&Cartan matrix, roots, $d$, $\dim\fg$, $h^\vee$&
tensor decomposition, projectors, $K,T,\pi$, PBW/Koszul, $\vartheta$, tower maps\\
$G_2$&$7$&four recorded channels and their Casimir scalars&
four projectors, $K,T,\pi$, PBW/Koszul, $\vartheta$, tower maps
\end{tabular}
\end{table}

\begin{table}[ht]
\centering
\small
\caption{Passage from a finite quadratic presentation to the completed Yangian.}
\label{tab:exceptional-yangian-ambient-registers}
\begin{tabular}{c|p{4.8cm}|p{5.4cm}}
\textsc{stage}&\textsc{object}&\textsc{map}\\ \hline
finite coefficients&$W_N(V_0)$ with its trace pairing&
$R_N(K)^\perp\xrightarrow{\vartheta_N}R_N(K^{-1})$\\
finite algebra&$A_N(K)=T(W_N)/\langle R_N(K)\rangle$&
$\pi_{\fg,N}\colon A_N(K)/\mathcal N_N\xrightarrow{\sim}Y_{\hbar,\le N}(\fg)$\\
normalization&$\mathcal R_{\fg,N}=\operatorname{Rees}_F(A_N/\mathcal N_N)$&
specialization at $t=1$ to $Y_{\hbar,\le N}$ and at $t=0$ to
$\operatorname{gr}_F Y_{\hbar,\le N}$\\
bar construction&$A_N^{\mathrm i}\xrightarrow{q_N}B(A_N)$&
$q_N$ is a quasi-isomorphism under diagonal Koszulness\\
completion&$\varprojlim_N A_N(K)$&
$p_{N+1,N}$ and weightwise Mittag--Leffler passage\\
chiralization&$\mathfrak C_X(A_N)$&
$\gamma_N\colon\mathfrak C_X(A_N^{!,\mathrm{quad}})\xrightarrow{\sim}
\mathfrak C_X(A_N)^{!,\mathrm{ch}}$
\end{tabular}
\end{table}

\begin{table}[ht]
\centering
\small
\caption{Five typed objects associated with an RTT algebra.}
\label{tab:exceptional-yangian-object-firewall}
\begin{tabular}{c|p{9.5cm}}
\textsc{symbol}&\textsc{definition}\\ \hline
$A_N$&$T(W_N)/\langle R_N(K)\rangle$\\
$B(A_N)$&$T^c(s^{-1}\bar A_N)$ with the multiplication differential\\
$A_N^{\mathrm i}$&the quadratic coalgebra on $s^{-1}W_N$ with
corelations determined by $R_N(K)$\\
$A_N^{!,\mathrm{quad}}$&$T(W_N^\vee)/\langle R_N(K)^\perp\rangle$\\
$Z_{\mathrm{ch}}^{\mathrm{der}}(\mathfrak C_X(A_N))$&the chiral
Hochschild cochain object of the chiralized algebra
\end{tabular}
\end{table}

\begin{remark}[Passage from root data to exceptional duality;
\ClaimStatusProvedHere]
\label{rem:exceptional-yangian-three-table-quarantine}
Tables~\ref{tab:exceptional-yangian-theorem-status},
\ref{tab:exceptional-yangian-ambient-registers}, and
\ref{tab:exceptional-yangian-object-firewall} display the complete
sequence of maps.  Cartan arithmetic chooses $\fg$ and $V_0$; projector
matrices and $K_{\fg,\hbar}$ define $R_N(K)$; $\pi_{\fg,N}$ identifies
the RTT algebra with the Yangian; $q_N$ and $\vartheta_{\fg,N}$ give
quadratic duality; the transition maps give the completion; $\gamma_N$
gives the chiral formulation.
\end{remark}
